\chapter{The frontier: where covering stops, and why}\label{ch:frontier}

With a calibrated exponent in hand we can ask sharp questions. How small can a
desert be made at a given target $Q$? How large can $\g$ be made under a given
digit budget? And---the question this chapter is really about---is the answer
determined by mathematics we could improve, or by a wall?

The answer turns out to be surprisingly clean, and surprisingly unlike what we
expected when we started. The difficulty of the whole subject is an
\emph{integrality gap}: fractionally, the covering problem is trivial. All of
the exponent, every last nat of it, is the price of rounding a fractional cover
to an integral one.

\section{The frontier, measured}

Table~\ref{tab:frontier} gives the achievable exponent as a function of the
digit budget, at fixed target $Q=10^{5}$: for each cap on the number of digits
of $\N$, we build the best cover we can and report $|U|$ and $E$.

\begin{table}[h]
\centering\small
\caption{Cover frontier at $Q\approx10^{5}$: the exponent achievable under a
digit cap. Fewer digits force a smaller modulus, hence fewer moduli, hence a
larger residual set \emph{and} a smaller $\log\N$ to divide by---the two effects
compound, which is why the frontier steepens so sharply below $130$
digits.}\label{tab:frontier}
\begin{tabular}{rrrr}
\toprule
Digits of $\N$ & Moduli & $|U|$ & $E$\\
\midrule
$195$ & $89$ & $692$ & $16.8$\\
$170$ & $79$ & $754$ & $20.6$\\
$150$ & $72$ & $802$ & $24.4$\\
$140$ & $68$ & $830$ & $26.8$\\
$130$ & $64$ & $862$ & $29.5$\\
$120$ & $60$ & $899$ & $32.8$\\
$100$ & $52$ & $973$ & $41.3$\\
\bottomrule
\end{tabular}
\end{table}

Table~\ref{tab:ladder} does the complementary experiment: fix the digit budget
at $199$ and raise the target. This is the budget game's own ladder, and it is
the table that governs Chapter~\ref{ch:outlook}'s planning.

\begin{table}[h]
\centering\small
\caption{The rung ladder at $199$ digits: exponent and candidate count as the
target $Q$ rises, using the best cover we can build at each target ($88$ moduli,
$\mathrm{boost}\approx11$). Fleet-days assume the measured
$7.9\times10^{5}$ candidates/s of a two-session GPU fleet.}\label{tab:ladder}
\begin{tabular}{rrrrr}
\toprule
$Q$ & $|U|$ & $E$ & $e^{E}$ candidates & fleet time\\
\midrule
$122\,000$ & $867$ & $20.83$ & $1.1\cdot10^{9}$ & $23$ minutes\\
$135\,000$ & $1\,004$ & $24.12$ & $3.0\cdot10^{10}$ & $11$ hours\\
$150\,000$ & $1\,138$ & $27.34$ & $7.5\cdot10^{11}$ & $11$ days\\
$175\,000$ & $1\,335$ & $32.07$ & $8.5\cdot10^{13}$ & $3.4$ years\\
$200\,000$ & $1\,569$ & $37.69$ & $2.4\cdot10^{16}$ & $950$ years\\
\bottomrule
\end{tabular}
\end{table}

The spacing is geometric---each rung costs $25$ to $280$ times the last---which
has a pleasant strategic consequence exploited in Chapter~\ref{ch:outlook}:
climbing \emph{every} rung costs only about $4\%$ more than jumping straight to
the top, and banks a certified record at each step.

\section{Cover size is an equilibrium, not a budget}\label{sec:equilibrium}

A natural mental model is that the cover is limited by the digit budget: use as
many moduli as $\log M\le\log\N$ allows. That model is wrong at record scale,
and correcting it explains why the frontier is so flat.

Consider adding one more modulus $r=461$ to the $88$-modulus cover behind our
budget-game work. It kills about $3.5$ residual offsets, worth
$-3.5\times\mathrm{boost}/\log\N=-0.084$ nats. But it also multiplies
$\mathrm{boost}$ by $461/460$, which raises the primality odds of every
\emph{surviving} residual complement, worth $+0.082$ nats. The net effect is
$+0.00$ nats: the cover is at a stationary point.

\begin{observation}[Boost--kills equilibrium]\label{obs:equilibrium}
At $B=10^{200}$ the cover behind our budget-game work spends $\log M\approx437$
of the $\approx461$ available nats, so the digit budget is not tight; the cover
stops growing because $dE=0$, not because it has run out of room. Adding moduli
past the equilibrium is neutral to slightly harmful.
\end{observation}

This is why leftover digit budget buys nothing, why the frontier in
Table~\ref{tab:frontier} is smooth rather than kinked, and why---as we will
see---only two things can actually move $E$: better rounding mathematics, or
more compute.

\section{The linear-programming relaxation is vacuous}\label{sec:lp}

The covering problem is a set-cover: choose one residue class per modulus so as
to cover as many prime offsets as possible. Set-cover problems come with a
standard relaxation---allow fractional choices---whose optimum lower-bounds the
integral one and whose dual prices certify optimality. A natural research
programme, proposed in our own earlier planning documents, was to solve that
relaxation and publish the duals as a certificate of how close greedy is to
optimal.

We solved it. The result kills the programme, and in doing so reframes the whole
subject.

\begin{proposition}[LP vacuousness]\label{prop:lp}
Let $\mathcal R$ be the $88$ moduli $3,5,\dots,457$ used at record scale and let
the universe be the odd primes below $Q$. In the relaxation where each modulus
$r$ may distribute a unit of weight over its residue classes and a prime is
covered to the extent that weight sits on its class, the optimum covers
$100.00\%$ of the primes---at both $Q=122\,000$ and $Q=200\,000$.
\end{proposition}
\begin{proof}[Verification]
Solved with HiGHS in about two seconds at each target. The reason is a mass
count: modulus $r$ covers a $1/(r-1)$ fraction of the odd primes in expectation,
and $\sum_{r\in\mathcal R}1/(r-1)\approx1.86$. With $186\%$ of the required
covering mass available, a fractional assignment can spread it to leave nothing
uncovered.
\end{proof}

So the LP dual certifies nothing whatsoever about integral covers, and the
$\S7.4$-style ``optimality certificate'' programme is dead as posed. What
replaces it is a much better statement:

\begin{corollary}[The exponent is a rounding gap]\label{cor:rounding}
Every nat of the failure exponent is the cost of integrality. Concretely, at
$Q=122\,000$:
\[
\underbrace{15.6\%}_{\text{naive rounding}}\;\longrightarrow\;
\underbrace{7.56\%}_{\text{adaptive greedy}}\;\longrightarrow\;
\underbrace{0\%}_{\text{fractional optimum}}
\]
uncovered primes, where the first figure is $e^{-1.86}$, the density left by
independent rounding at the available mass. Each absolute percentage point of
residual removed is worth about $2.7$ nats---a factor of $15$ in compute.
\end{corollary}

Greedy therefore sits almost exactly halfway, in halvings, between doing nothing
clever and perfection. The rest of this chapter asks whether the remaining half
can be taken.

\section{Attacking the rounding gap}\label{sec:door1}

We fix a concrete benchmark: the $88$ moduli above, the $11\,474$ odd primes
below $Q=122\,000$, and the greedy/coordinate-descent incumbent with
$|U|=867$. One prime saved is worth $0.0240$ nats. The question is how much of
the gap to $0$ is reachable.

\subsection{How much is there to win?}

Before optimising, it is worth knowing the size of the prize. Sampling the space
of assignments uniformly (choosing each $a_r$ at random) gives
$|U|$ with mean $\mu=1377$ and standard deviation $\sigma=21.0$. Two things are
notable. First, $\sigma$ is \emph{below} the independence prediction of $34.8$:
the choices are negatively correlated, so the deep tail is thinner than a naive
model suggests. Second, the design space has entropy
$\ln\prod_r(r-1)=435.7$ nats, so a first-moment (annealed) calculation puts the
best assignment in the ensemble at roughly
\[
u^{*}\;\approx\;\mu-\sqrt{2\ln\Omega}\,\sigma\;\approx\;758 ,
\]
against greedy's $867$. Greedy already sits at $z=24.3$ of the $29.5$ standard
deviations theoretically available. The honest prize is therefore about
$2.6$ nats---a factor of $14$---and only if the annealed threshold is
achievable, which for a clustered landscape it typically is not.

\subsection{Exact local optimisation}

We attacked the benchmark with large-neighbourhood search: freeze all but a
window of $7$--$15$ moduli and re-solve that window \emph{exactly} with
CP-SAT~\cite{CPSAT} over the primes left uncovered by the frozen part, then
rotate windows. Windows were chosen three ways---by tightness (the moduli whose
current class beats its best alternative by the smallest margin), randomly among
tight moduli, and uniformly---and the incumbent was enforced as a hard floor so
the solver could only match or improve.

Across $41$ windows, with time caps up to $300$ seconds, the incumbent was
matched every time and never beaten. To turn ``not beaten'' into ``provably
optimal at this scale'' we added an LP certificate per window: solve the
window's own relaxation, and if $\lfloor\mathrm{LP}\rfloor$ equals the
incumbent's coverage then no integral assignment of that window can do better.
This certified $21$ of $22$ windows where it was computed. The single exception
was a window of the smallest moduli---where, exactly as
Proposition~\ref{prop:lp} predicts, the LP is loose.

We then went further and solved the \emph{entire large-modulus endgame} at once:
first the $30$ largest moduli simultaneously, then the $44$ largest---half the
cover in a single exact subproblem. Both leave $|U|=867$.

\subsection{Global methods}

Local exactness having failed, we tried methods that can cross basins.

\emph{Belief propagation with decimation.} We wrote the assignment problem as a
factor graph---one variable per modulus with domain $\mathbb{Z}/r$, one factor
per prime contributing weight $e^{-\beta}$ when uncovered---and ran sum-product
message passing, then repeatedly froze the most polarised modulus and re-ran
(the standard decimation scheme~\cite{MezardMontanari}). Validation first: on a
small instance where exhaustive enumeration is possible, BP tracks the exact
$E[|U|]$ to $0.7\%$ in the soft regime, and BP-guided decimation finds the exact
optimum. On the real instance it returns $874$ at $\beta=1.5$ (and worse at
higher $\beta$, where early freezes over-commit).

That is seven primes short of greedy, and it is the second-best construction we
know for this instance---but it is not an improvement.

\emph{Basin structure.} The most informative experiment was the cheapest. Take a
solution produced by random-start coordinate ascent, which lands at
$|U|\approx968$, and apply the same exact-window machinery to it: twenty
windows recover \emph{none} of its $101$-prime gap to greedy. Independent ascent
solutions agree with each other on only $2.4\%$ of their class choices---chance
level---so there is no backbone to exploit, and crossover-style population
methods have nothing to graft.

\begin{observation}[The landscape]\label{obs:landscape}
The assignment space is a field of mutually distant, locally-exact-stable
basins spanning at least $|U|\in[867,980]$. Exact local search cannot move
between them at all; what distinguishes greedy is not the neighbourhood it
explores but the \emph{sequentially adaptive} way it constructs a solution,
which lands it about $100$ primes deeper than a typical basin.
\end{observation}

\subsection{Verdict, and one methodological caution}

The greedy plateau survived every method we brought: $41$ exact windows (mostly
LP-certified), a half-cover exact block, BP decimation, and basin probing. We
state the conclusion carefully, because the temptation to over-claim here is
strong.

\begin{observation}[Door-1 verdict]\label{obs:door1}
At $Q=122\,000$ with these moduli, $|U|=867$ is at or very near the quenched
optimum. The annealed estimate of $758$ is an \emph{ensemble-existence}
heuristic, not a constructive promise, and the clustering visible in
Observation~\ref{obs:landscape} is exactly the situation in which annealed and
quenched optima diverge. The realistic prize is therefore $0$ to $2.6$ nats,
probably nearer the bottom of that range.
\end{observation}

The caution concerns mean-field methods. Naive belief propagation, run to low
temperature, predicts $E[|U|]\to35$ at $\beta=4$---heading for the fractional
fantasy of Proposition~\ref{prop:lp} rather than the achievable $867$. The
replica-symmetric approximation inherits precisely the LP's blindness to
integrality. Anyone using message passing to estimate a threshold here needs the
one-step replica-symmetry-breaking treatment with complexity bookkeeping; we
flag this as the honest remaining theory item rather than quoting a number we do
not trust.

\section{What the frontier costs in money}

Putting Table~\ref{tab:frontier} together with measured throughput gives the
price list in Table~\ref{tab:menu}. It is worth stating in these terms because
the strategic question in Chapter~\ref{ch:outlook} is precisely which of these
rungs is worth buying.

\begin{table}[h]
\centering\small
\caption{Cost of the threshold game at $Q=10^{5}$ by digit target, at
$8.5\times10^{6}$ tests/s per A100 and cloud pricing current at the time of
writing.}\label{tab:menu}
\begin{tabular}{rrrr}
\toprule
Digits & $E$ & A100-years & approx.\ cost\\
\midrule
$140$ & $28.1$ & $0.10$ & \$1\,000\\
$135$ & $29.7$ & $0.49$ & \$5\,000\\
$130$ & $31.2$ & $2.3$ & \$24\,000\\
$125$ & $32.9$ & $12$ & \$125\,000\\
$120$ & $34.7$ & $73$ & \$770\,000\\
$100$ & $44.1$ & $8.9\cdot10^{5}$ & \$9.4\,billion\\
\bottomrule
\end{tabular}
\end{table}

The table makes the shape of the problem vivid. The frontier is fundable down to
about $130$ digits and institutional down to $120$; sub-$100$ is seven orders of
magnitude beyond plausible compute. Whether that last wall can be attacked by
mathematics rather than money is the subject of the next chapter.
